\section{The ideal structure of \texorpdfstring{$\Squad$}{Squad}}\label{sec:ideal}

The designers' complexity estimate assumes that the $3(\ell+1)n$ quadratics of $\Squad$ behave like a
semi-regular system. This section shows that they do not: two thirds of them are redundant by
construction, the degree-3 Macaulay matrix has $8n$ syzygies per S-box, and the first-fall degree is
$4$ with a fall count that depends only on the multiplicative class of $e$ modulo $n$.

\subsection{Two thirds of \texorpdfstring{$\Squad$}{Squad} is an XL step}

%% ledger: F
\begin{proposition}\label{prop:ideal-P1}
\Proved\ In the Boolean ring, $\langle P_1,P_2,P_3\rangle=\langle P_1\rangle$; more precisely each
coordinate of $P_2$ and of $P_3$ lies in the $\F_2$-span of $\{z\cdot p : z \text{ a variable},\ p
\text{ a coordinate of } P_1\}$. Hence $\Squad$ is exactly the first XL step applied to the $P_1$
system, and $d_{\mathrm{solv}}(\Squad)\ge d_{\mathrm{solv}}(P_1)-1$.
\end{proposition}
\begin{proof}
$P_2=xP_1$ and $P_3=yP_1$ as elements of $\Fq[x,y]$. For any $\gamma$,
$\Tr(\gamma\, xP_1)=\sum_j x_j\,\Tr(\gamma\beta_j P_1)$ where $x=\sum_j \beta_j x_j$ is the expansion
of $x$ in the chosen basis; each $\Tr(\gamma\beta_j P_1)$ is a coordinate of $P_1$. The same argument
applies to $P_3$. The converse inclusion is trivial.
\end{proof}

\subsection{Syzygies at degree three}

%% ledger: G
\begin{theorem}[$8n$ syzygies]\label{thm:syz}
For one S-box, $\syz_3\bigl(\{P_1,P_2,P_3\}\bigr)=8n$. Of these, $3n$ are proved
(Proposition~\ref{prop:trivial-syz}); the remaining $5n$ are \Verified{$n=11..16$, every admissible
$e$} --- $36$ runs, no exception --- and \Conj\ in general.
\end{theorem}

%% ledger: G
\begin{proposition}[the $3n$ structural syzygies]\label{prop:trivial-syz}
\Proved\ The relations $x\cdot P_1=P_2$, $y\cdot P_1=P_3$ and $y\cdot P_2=x\cdot P_3$ each contribute
$n$ independent linear dependencies among the rows of $\Mac_3$.
\end{proposition}
\begin{proof}
As in Proposition~\ref{prop:ideal-P1}, every coordinate of $P_2$ is an $\F_2$-combination of the rows
$x_j\cdot(\text{coordinate of }P_1)$; the $n$ coordinates of $P_2$ therefore give $n$ dependencies
between the degree-$0$-multiplier rows of $P_2$ and the degree-$1$-multiplier rows of $P_1$, and they
are independent because the coordinates of $P_2$ are. The same for $P_3$, and again for
$yP_2=xyP_1=xP_3$.
\end{proof}

The mechanism behind the remaining $5n$ is Frobenius-orbit cancellation in the sense of
\eqref{eq:trace-frobenius}. \todo{Write the $5n$ generators explicitly. The batch-2 search
(Section~\ref{sec:falls}) shows that the analogous question at $D=4$ has no small structured answer,
so this is expected to be delicate.}

\subsection{First-fall degree four}\label{sec:falls}

%% ledger: H
\begin{theorem}[first-fall degree and fall count]\label{thm:firstfall}
\Verified{$n=11..16$, every admissible $e$; $36$ runs} For one S-box the first-fall degree is $4$,
and the number of new cubics is
\[
  \mathrm{falls}_4 = 2n+\Delta(e),
\]
where $\Delta$ depends only on the unordered pair $\{e,\,e^{-1}\bmod n\}$ and is given by
Table~\ref{tab:classes}. In particular $\Delta(e)=0$ --- i.e.\ exactly $2n$ falls --- for every
exponent of every recommended parameter set.
\end{theorem}

%% ledger: I
\begin{table}[t]\centering
\caption{The measured fall count as a function of the class of $e$. Source:
\texttt{results/E03\_summary.md} (36 runs, $n=11..16$). The rule reproduces every measured count;
it is a fit to $n\le16$, hence \Conj.}\label{tab:classes}
\begin{tabular}{ll}
\toprule
class of $e$ (as a pair $\{e,e^{-1}\}$) & $\mathrm{falls}_4$\\
\midrule
$\{2,\,2^{-1}\}$ & $7n$\\
$\{-2,\,(-2)^{-1}\}$ & $4n$\\
$\{1+n/2\}$ ($n$ even, self-inverse) & $17n/2$\\
$\{-3,\,(-3)^{-1}\}$ & $2n+1$\\
$\{n/2-1\}$ ($n$ even, self-inverse) & $2n+1$\\
generic & $2n$\\
\bottomrule
\end{tabular}
\end{table}

%% ledger: I2
\begin{observation}[the sign matters]\label{obs:sign}
The extra fall of the $2n+1$ class appears exactly when $3e\equiv-1$ or $3e^{-1}\equiv-1 \pmod n$, or
$e=n/2-1$. The congruence $3e\equiv+1$ does \emph{not} produce it: $(n,e)=(11,4),(13,9),(14,5),(16,11)$
all satisfy it and all give exactly $2n$. Every other family tested --- $k\cdot e\equiv\pm1,\pm2$ for
$k\le6$, for $e$ and for $e^{-1}$; $e$ against $n/2$; $e^2\equiv\pm1$ --- fails to separate the two
classes.
\end{observation}

%% ledger: H2, H3
\begin{theorem}[no cubic generators outside the quadratic ideal]\label{thm:outside}
\Verified{$n=11..16$, every admissible $e$ with $e\ne2$} The degree-$4$ Macaulay matrix exhausts the
cubic part of the vanishing ideal of the S-box graph: $\dim I(V_e)_{\le3}=\mathrm{low}_4$. The
exception is $e=2$, where exactly $2n$ cubics of $I(V_e)$ are not reached at $D=4$ ($n=11,13,15$).
Curiously $\mathrm{falls}_4$ is symmetric under $e\leftrightarrow e^{-1}$ but this defect is not:
$e=2^{-1}$ has no unreachable cubics. \Open
\end{theorem}

\subsection{What generates the falls}

The natural guess is that the $2n$ new cubics come from a small, structured family of field
identities, in the same way that $3n$ of the syzygies come from $P_2=xP_1$ and $P_3=yP_1$. Because a
squarefree Boolean monomial of degree $\le2$ is itself a Boolean-degree-$\le2$ function of $(x,y)$,
every row of $\Mac_4$ is of the form $\Tr(\gamma\,M\,P_k)$ for a \emph{field} monomial $M=x^ay^b$ with
$\wt(a)+\wt(b)\le2$; and the degree-$4$ part of $\Tr(\gamma\,x^{2^i}y^{2^j}P_1)$ is
$\Tr\bigl(\gamma\,x^{2^i+1}y^{2^j+1}\bigr)$, which is homogeneous of degree $\wt(2^i+1)+\wt(2^j+1)$.
A fall is therefore a solution of a pure Frobenius-cancellation condition, and the question is which
multiplier families are needed to realise one.

\begin{observation}[the natural ansatz is refuted]\label{obs:ansatz}
\Verified{$(n,e)=(11,3),(13,3)$} The $\Fq$-module generated by
$\{x^{2^i}y^{2^j}P_1\}\cup\{x^{2^k}P_2\}\cup\{y^{2^k}P_3\}$ does produce degree-$4$ cancellations,
but every cubic it yields already lies in $\Mac_3$: it explains \emph{none} of the $2n$ new cubics.
Moreover none of the $18$ multiplier families
$\{1,x^{2^i},y^{2^j},x^{2^i}y^{2^j},x^{2^i+2^j},y^{2^i+2^j}\}\times\{P_1,P_2,P_3\}$ produces a single
new cubic on its own.
\end{observation}

Dropping one family at a time from the full row space isolates the two that matter.
%% ledger: R
\begin{observation}[two indispensable families]\label{obs:two-families}
\Verified{$(n,e)=(11,3)$; all 18 families measured} Removing any of 16 families leaves the fall count
at $2n$; removing $x^{2^i}y^{2^j}\cdot P_2$, or removing $y^{2^i+2^j}\cdot P_3$, leaves exactly $n$.
Hence $\mathrm{falls}_4=2n$ splits as $n+n$, and both blocks need \emph{weight-two} multipliers on
$P_2$, $P_3$ --- not the weight-one ones of Observation~\ref{obs:ansatz}.
\end{observation}

The reason is visible in the degree-4 count. Writing every generator as a multiple of $P_1$:
$x^{2^k}P_2=x^{2^k+2}y+x^{2^k+1+2^e}+x^{2^k+1}$ has \emph{no} degree-4 part at all (Boolean degree
$\le3$), so it lies in $\Mac_3$ and can contribute nothing; $x^{2^i}y^{2^j}P_1$ has exactly
\emph{one} degree-4 monomial, so its only cancellations are the $\lfloor(n-1)/2\rfloor$ diagonal
Frobenius coincidences, and those were measured to land in $\Mac_3$ as well; whereas
\[
  x^{2^i}y^{2^j}P_2=x^{2^i+2}y^{2^j+1}+\mathbf{x^{2^i+1+2^e}y^{2^j}}+x^{2^i+1}y^{2^j},\qquad
  y^{2^i+2^j}P_3=xy^{2^i+2^j+2}+\mathbf{x^{2^e}y^{2^i+2^j+1}}+y^{2^i+2^j+1}
\]
each carry \emph{two} degree-4 monomials, the second of which involves $2^e$. Cancellation can
therefore occur inside a single generator, and the condition sees the exponent --- which is how
$\mathrm{falls}_4$ can depend on the class of $e$ at all. A single-monomial family could not.

\begin{remark}[an obstruction to the obvious shape]\label{rem:obstruction}
$\mathrm{falls}_4=2n$ invites the description ``two field cubics, $n$ coordinates each''. That cannot
be the general shape: the classes $\{-3,(-3)^{-1}\}$ and $\{n/2-1\}$ give $2n+1$
(Table~\ref{tab:classes}), and $2n+1$ is not a multiple of $n$. Any correct statement of
Theorem~\ref{thm:firstfall} must produce a lone extra element there, not a third $n$-block.
\end{remark}

\subsection{Across S-boxes: no cubic relations, and the direct sum}

%% ledger: K, K2
\begin{theorem}[no cross-S-box cubic]\label{thm:cross}
\Verified{$n=14,16$} Let $t_0,t_1$ be the outputs of two first-layer S-boxes driven by the same $pt$.
The space of Boolean cubics vanishing on $\{(t_0,t_1)\}$ consists exactly of the univariate image
criteria of Proposition~\ref{prop:image}; its dimension is $\#\{i:\ n-m_i\le3\}$. There is no cubic
relation involving both $t_0$ and $t_1$.
\end{theorem}
\begin{remark}
At $n=14$ with $e=(3,5)$ the criteria have Boolean degree $9$ and $11$ and the dimension is $0$; at
$n=16$ with the same exponents, $e_1=5$ has $m_1=13$ so the criterion has degree $n-m_1=3$ and the
dimension is $1$. Reading that $1$ as a cross relation would be a mistake: the relation is supported
on the $t_1$ variables alone. This is the same quantity $n-m$ that Corollary~\ref{cor:degrees}
identifies as a design criterion.
\end{remark}

%% ledger: M, M2, M3
\begin{theorem}[direct sum at low degree]\label{thm:directsum}
For the full system, $\syz_3(\Squad)=(\ell+1)\cdot 8n$ \Verified{$n=7,9,11$; $\ell=1,2$}. At $D=4$,
for the two-S-box first layer the falls are the exact direct sum of the per-S-box falls
(\Verified{$n=11$}: $44=22+22$ with zero cross excess), and ideal-theoretically
$\dim I(V)_{\le3}-\operatorname{rank}\Mac_3=4n$ exactly \Verified{$n=14$}.
\end{theorem}

%% ledger: M4
\begin{remark}[why the full three-S-box system is not measured here]\label{rem:e04}
For $\Squad$ itself the degree-$4$ row space saturates the degree-$\le3$ space at every size
reachable without a dedicated linear-algebra backend ($0.9936$--$0.9997$ of $\binom{N}{\le3}$ for
$n\le11$), and a random quadratic system of the same shape saturates identically. The fall count
there measures the monomial basis, not the algebra (Remark~\ref{rem:saturation}), so no conclusion is
drawn from it. \todo{Server batch: $n=13,15$, where the matrix is no longer saturated.}
\end{remark}

\subsection{The linear layer is not the source of any of this}

%% ledger: P
\begin{observation}[matrix distribution is irrelevant]\label{obs:distribution}
\Verified{$n=9,11$; two seeds} Replacing the specification's unit-LU matrices by uniform random
elements of $GL_n(\F_2)$, or by random invertible matrices of the same density, changes neither the
degree-3 rank and syzygy count nor the degree-4 fall count of the first-layer system. At $n=11$ all
three distributions give rank $2068$, $\syz_3=176=2\cdot 8n$ and $\mathrm{falls}_4=44=4n$.
\end{observation}

Everything in this section is therefore a property of the S-box, not of the key schedule.
